\documentclass[11pt,a4paper]{article}
\usepackage[utf8]{inputenc}
\usepackage[T1]{fontenc}
\usepackage{lmodern}
\usepackage[spanish]{babel}
\usepackage[margin=2.1cm]{geometry}
\usepackage{tabularx}
\usepackage{fancyvrb}
\usepackage{microtype}
\usepackage{xurl}
\usepackage{hyperref}
\hypersetup{colorlinks=true,urlcolor=blue,linkcolor=black,
 pdftitle={Laboratorio 10 - Captura el Router},
 pdfauthor={Javier Rivera y Erick Mejia}}
\DefineVerbatimEnvironment{comandos}{Verbatim}{fontsize=\footnotesize,frame=single}
\setlength{\parindent}{0pt}
\setlength{\parskip}{6pt}
\newcommand{\carneJose}{18007794}
\newcommand{\carneErick}{20006104}
\begin{document}

\begin{center}
{\Large\textbf{Laboratorio 10: Captura el Router}}\\[4pt]
{\large Ciencias de la Computación VIII}\\
Redes y Protocolos\\[12pt]
Javier Rivera --- Carné: \carneJose\\[7pt]
Erick Mejia --- Carné: \carneErick\\[10pt]
20 de octubre de 2026
\end{center}

\section{Objetivo y alcance de la entrega}
Preparar la captura de tráfico del router WEP dispuesto por el docente y la
recuperación estadística de su clave, sin diccionarios. Se documentan comandos,
automatización mínima, evidencia reproducible y condiciones adversarias que
pueden reducir la calidad de captura.

\textbf{Estado de este informe: preparación previa a la competencia.} Se ejecutó
una recuperación real sobre una captura pública de Aircrack-ng y se verificó el
ayudante local. El router y el adaptador físico se probarán en clase. No se
presentan resultados de captura en vivo ni acciones contra compañeros como si
ya hubieran ocurrido. El enunciado establece evaluación individual; la portada
incluye los dos nombres solicitados para esta preparación.

\subsection*{Fundamento técnico}
WEP utiliza RC4 con una clave secreta compartida y un vector de inicialización
(IV) de 24 bits que viaja con cada trama. El espacio limitado de IVs y las
debilidades estadísticas de la construcción permiten inferir la clave a partir
de tráfico suficiente. PTW aprovecha correlaciones del flujo RC4; no necesita
adivinar palabras de una contraseña. Una clave WEP de ``64 bits'' incluye
40 bits secretos más el IV de 24 bits; la de ``128 bits'' incluye 104 bits
secretos. El trabajo original de PTW sustenta esta recuperación estadística [1].

Aircrack-ng usa PTW por defecto para WEP y permite seleccionar WEP con
\texttt{-a 1}. Se capturan paquetes completos para conservar el material útil,
incluido ARP; no se usa \texttt{-w} ni una captura limitada con \texttt{--ivs}.
No se confunde esta práctica con la obtención de un handshake WPA [2].

\subsection*{Entregables preparados}
\path{informe.pdf}; ayudante \path{lab10.py}; doce pruebas en
\path{test_lab10.py}; muestra pública \path{muestras/ptw.cap}; registros
locales en \path{evidencias/}; plantilla de medición para clase. Este documento
debe subirse a GES antes de la evaluación, después de completar los carnés.

\newpage
\section{Herramientas, equipo y preparación}
\begin{tabularx}{\linewidth}{|l|X|}\hline
Herramienta & Función en el laboratorio \\\hline
\texttt{airmon-ng} & Detectar interfaces y activar/restaurar modo monitor. \\\hline
\texttt{airodump-ng} & Guardar tráfico 802.11 y observar IVs/contadores. \\\hline
\texttt{aireplay-ng} & Prueba activa de inyección y reinyección ARP opcionales. \\\hline
\texttt{aircrack-ng} & Recuperación offline de la clave WEP. \\\hline
\texttt{iw}, Python 3 & Revisar hardware e interfaces; ejecutar el ayudante. \\\hline
Wireshark (opcional) & Examinar tramas, reintentos y gestión 802.11. \\\hline
\end{tabularx}

En Ubuntu/Debian se instalan las cuatro herramientas Aircrack desde un paquete:
\begin{comandos}
sudo apt update
sudo apt install aircrack-ng iw python3
# Opcional para analizar visualmente las tramas:
sudo apt install wireshark
\end{comandos}

Se requiere una \textbf{interfaz Wi-Fi física compatible con modo monitor};
para la ruta activa también debe admitir inyección. En una VM hay que asignarle
el adaptador USB directamente. Una interfaz Ethernet virtual con Internet no
proporciona acceso a tramas 802.11 del aire.

El entorno de preparación fue Ubuntu sobre WSL2, kernel
\texttt{6.18.54.1-microsoft-standard-WSL2}. Al verificar estaban instalados
Aircrack-ng 1.7 e \texttt{iw}, pero no se detectó interfaz Wi-Fi. Por ello aquí
se puede analizar una captura guardada, pero no se ha validado radio.
USB/IP permite conectar dispositivos USB a WSL2; no garantiza que el chipset y
su controlador admitan monitor e inyección [7].

\subsection*{Comprobaciones que sí pueden hacerse sin el aula}
Desde el directorio del laboratorio:
\begin{comandos}
python3 lab10.py comprobar
python3 -m unittest -v test_lab10.py
sha256sum -c evidencias/muestra.sha256
aircrack-ng -a 1 -b 00:12:BF:12:32:29 muestras/ptw.cap
\end{comandos}
La muestra se incluye para practicar sin Internet. Su clave pública esperada
es \texttt{1F:1F:1F:1F:1F}, según el README oficial [6]. No es la del aula.
\texttt{comprobar} devuelve 1 si faltan herramientas o no detecta Wi-Fi; ese
resultado es esperado en este WSL2. Las pruebas locales no necesitan privilegios
de administrador; la prueba de recuperación real se omite explícitamente si
Aircrack-ng no está instalado.

\newpage
\section{Procedimiento de captura y recuperación en clase}
\subsection*{1. Identificar únicamente el objetivo autorizado}
Solicitar al docente SSID, BSSID y canal. No basta con el nombre de la red: el
BSSID identifica el AP. Los valores siguientes son ejemplos que se deben
\textbf{sustituir por los reales}; el BSSID ilustrativo no corresponde al aula.
\begin{comandos}
IFACE=wlan0
MON=wlan0mon
BSSID=02:11:22:33:44:55
CANAL=6
sudo airmon-ng
iw list
sudo airmon-ng start "$IFACE" "$CANAL"
iw dev
\end{comandos}
En \texttt{iw list} se busca soporte de monitor; luego se comprueba el tipo de
interfaz en \texttt{iw dev}. Ajustar \texttt{MON} al nombre que realmente crea
\texttt{airmon-ng} [3]. Si el docente no facilita el canal, observar pasivamente
con \texttt{sudo airodump-ng "\$MON"}, identificar el AP autorizado y detener
esa observación antes de iniciar la captura con canal fijo.

\subsection*{2. Capturar paquetes completos}
Mantener fija la sintonía evita perder tiempo escuchando otros canales.
\begin{comandos}
mkdir -p capturas
sudo airodump-ng --bssid "$BSSID" --channel "$CANAL" \
  --encrypt WEP --write capturas/clase \
  --output-format pcap,csv "$MON"
\end{comandos}
Confirmar que el AP muestra WEP y que aumenta \texttt{\#Data}. Se generan
\texttt{clase-01.cap} y \texttt{clase-01.csv}; si el prefijo ya existe, el
sufijo puede ser otro. Conservar la ruta efectivamente creada [4].

\subsection*{3. Recuperar la clave en otra terminal}
Volver a definir el BSSID real en esa terminal y seleccionar la captura creada:
\begin{comandos}
BSSID=02:11:22:33:44:55
aircrack-ng -a 1 -b "$BSSID" \
  -l capturas/clave-clase.txt capturas/clase-01.cap
\end{comandos}
Se puede analizar mientras continúa la captura. Un resultado insuficiente
exige más tráfico útil o revisar la captura, no un diccionario. La cantidad de
IVs y el tiempo de éxito dependen de las condiciones; no se promete un umbral.
Registrar \texttt{KEY FOUND!}, el BSSID y el archivo utilizado, y confirmar el
resultado con el docente [2]. La salida hexadecimal es una clave WEP; no tiene
por qué corresponder a una frase legible.

\newpage
\section{Automatización, ruta activa opcional y restauración}
\subsection*{Ayudante desarrollado}
\texttt{lab10.py} automatiza el ciclo completo: diagnóstico (\texttt{comprobar}),
gestión del modo monitor (\texttt{monitor}), escaneo de APs WEP (\texttt{escanear}),
captura dirigida (\texttt{capturar}), generación de IVs (\texttt{inyectar}, con modos
\texttt{arpreplay}, \texttt{chopchop}, \texttt{fragment} y \texttt{fakeauth}),
desautenticación de clientes del AP (\texttt{deauth}), aislamiento de un host por
ARP spoofing (\texttt{bloquear}) y recuperación offline (\texttt{recuperar}).
Valida BSSID, canal, interfaz y direcciones IPv4, ejecuta argumentos sin un
intérprete de shell, conserva los errores de la herramienta y rechaza una salida
de clave ya existente. No implementa criptografía propia; todos los subcomandos
activos aceptan \texttt{--simular}, que valida y muestra el comando sin ejecutarlo
ni transmitir. Equivalentes a los pasos manuales anteriores:
\begin{comandos}
# Se pueden revisar los argumentos sin radio ni privilegios:
python3 lab10.py capturar --interfaz "$MON" \
  --bssid "$BSSID" --canal "$CANAL" --simular
# Terminal 1: captura real, con interfaz monitor ya preparada:
sudo python3 lab10.py capturar --interfaz "$MON" \
  --bssid "$BSSID" --canal "$CANAL" --salida capturas/clase
# Terminal 2: definir BSSID de nuevo y usar la captura real:
python3 lab10.py recuperar capturas/clase-01.cap \
  --bssid "$BSSID" --clave capturas/clave-clase.txt
\end{comandos}

\subsection*{Si falta tráfico útil: actividad dirigida al AP del laboratorio}
Esta ruta es activa, separada de la captura pasiva. Solo se usará dentro del
alcance establecido por el docente. La prueba de inyección envía probes,
incluidos probes broadcast aun al indicar un AP; un resultado de cero respuestas
no demuestra por sí solo incompatibilidad del adaptador [5].
\begin{comandos}
SSID='SSID indicado por el docente'
sudo aireplay-ng --test -e "$SSID" -a "$BSSID" "$MON"
\end{comandos}
Si el docente permite reinyección ARP en el AP de pruebas, se puede usar la MAC
de un \textbf{cliente propio o expresamente habilitado y ya asociado}. Mantener
\texttt{airodump-ng} capturando; detener la reinyección con Ctrl+C cuando proceda.
\begin{comandos}
CLIENTE=02:AA:BB:CC:DD:EE  # Sustituir por el cliente habilitado
sudo aireplay-ng --arpreplay -b "$BSSID" -h "$CLIENTE" "$MON"
\end{comandos}
La herramienta necesita observar un ARP válido: repetirlo induce al AP a
producir nuevas tramas con IVs nuevos. Sin cliente asociado, \texttt{inyectar
--modo fragment} o \texttt{chopchop} obtiene keystream para forjar un ARP, de modo
que la generación de IVs deja de depender del tráfico ajeno [8]. La desautenticación
(\texttt{deauth}) y el aislamiento por ARP (\texttt{bloquear}) sí están
automatizados y se reservan al entorno aislado autorizado por el docente, según se
detalla en la sección siguiente.

\subsection*{Restaurar el adaptador}
Detener captura y toda prueba activa con Ctrl+C; después ejecutar
\texttt{sudo airmon-ng stop "\$MON"}. Si se desactivó previamente la gestión
del adaptador con NetworkManager, reactivarla. No se ejecuta
\texttt{airmon-ng check kill} como paso automático: puede interrumpir la propia
conectividad. Ante cambios de canal, revisar primero
\texttt{sudo airmon-ng check} y la gestión del adaptador [3].

\newpage
\section{Condiciones adversarias y ventaja competitiva}
La sección adversarial del enunciado pide investigar congestión, pérdida y
disponibilidad. El docente precisó el alcance: la competencia se desarrolla sobre
una \textbf{red inalámbrica aislada} en la que solo participan los integrantes del
laboratorio, y autorizó expresamente acciones que degraden la eficiencia de captura
de los rivales. Dentro de ese alcance se desarrollaron y validaron en
\texttt{--simular} las acciones descritas abajo; su ejecución con radio se realiza
en el aula, sobre el AP de la práctica y nunca fuera de esa red aislada.

\begin{tabularx}{\linewidth}{|p{3cm}|X|X|}\hline
Condición investigada & Efecto esperado sobre la captura & Cómo observarla \\\hline
Congestión del canal por exceso de tráfico & Más contención y retransmisiones;
puede bajar la tasa de IVs únicos recibidos. Más paquetes totales no implica
siempre más información criptográfica. & Comparar IVs/s, RXQ y Lost con canal
fijo; revisar reintentos en Wireshark. \\\hline
Interferencia, obstáculos o mala ubicación & Menor recepción, tramas perdidas;
pueden verse beacons sin datos suficientes. & Comparar PWR y pérdida al mover
solo el propio receptor; el ruido no siempre es visible en una captura. \\\hline
Abuso de desautenticación o desasociación & Desconecta clientes y puede reducir
el tráfico útil. El receptor en monitor puede seguir viendo gestión: no
significa que desaparezcan todas las tramas. & Buscar gestión repetida y cambios
en estaciones; correlacionar con caída de IVs/s. \\\hline
Interferencia deliberada de radio & Puede degradar disponibilidad a más de un AP
y también perjudicar al propio capturador. & Una captura 802.11 por sí sola no
confirma interferencia no 802.11; puede requerir instrumentación de espectro. \\\hline
\end{tabularx}

Estas relaciones son hipótesis de diagnóstico, no resultados medidos en el aula.
\texttt{RXQ} resume recepción reciente; \texttt{Lost} estima huecos de secuencia
por estación; \texttt{\#Data} cuenta IVs únicos para WEP. Ninguno identifica por
sí solo al causante [4]. Wireshark permite filtrar reintentos o gestión del AP
autorizado sin transmitir [9]:
\begin{comandos}
wlan.bssid == 02:11:22:33:44:55 && wlan.fc.retry == 1
wlan.bssid == 02:11:22:33:44:55 &&
  (wlan.fc.type_subtype == 0x0c || wlan.fc.type_subtype == 0x0a)
\end{comandos}

\subsection*{Acciones desarrolladas para obtener ventaja}
\textbf{Autosuficiencia de IVs.} \texttt{fakeauth} asocia el adaptador al AP;
\texttt{fragment} o \texttt{chopchop} obtienen keystream aun con la red casi vacía
y \texttt{arpreplay} reinyecta un ARP forjado para producir miles de IVs sin
depender del tráfico de otros. Así la recuperación no se estanca por falta de
clientes.
\textbf{Degradación del rival.} \texttt{deauth} desautentica a los clientes del AP:
quien captura de forma pasiva se queda sin IVs, mientras el atacante sigue generando
los suyos por inyección (que no requiere clientes). \texttt{bloquear} realiza ARP
spoofing con sockets raw, sin dependencias externas, para aislar a un host tras
recuperar la clave; al no habilitar reenvío de IP su tráfico se descarta, y Ctrl+C
restaura la caché ARP. Ambos se ejecutan únicamente en la red aislada autorizada.
\textbf{Física.} Ubicar el propio adaptador donde mejore la recepción, con el cable
USB estable.
\textbf{Mitigación (lo que cierra estos ataques).} Todas funcionan porque WEP carece
de integridad y de protección de tramas de gestión: migrar a WPA2/WPA3 con 802.11w
neutraliza tanto la recuperación de clave como la desautenticación, y la inspección
dinámica de ARP junto con entradas estáticas detiene el \textit{spoofing}. No se
comparten procedimientos durante la prueba.

\newpage
\section{Evidencia técnica obtenida y pendientes}
\subsection*{Prueba offline ejecutada el 20 de octubre de 2026}
Se descargó \texttt{ptw.cap} desde el servidor oficial [6]. Se ejecutó el
ayudante con el BSSID de la muestra y se guardó la salida completa en
\texttt{evidencias/recuperacion\_muestra.txt}; se eliminaron únicamente secuencias
ANSI de presentación del terminal. El registro contiene comando, fecha y código
de salida. Extracto del resultado real:
\begin{comandos}
python3 lab10.py recuperar muestras/ptw.cap \
  --bssid 00:12:BF:12:32:29 --clave evidencias/clave_muestra.txt
Read 65282 packets.
Starting PTW attack with 30566 ivs.
KEY FOUND! [ 1F:1F:1F:1F:1F ]
Decrypted correctly: 100%
Codigo de salida: 0
\end{comandos}
El resultado coincide con la clave publicada. No se proporcionó esa clave al
programa: solo se comparó después. El archivo de salida contiene
\texttt{1F1F1F1F1F}. El SHA-256 de la captura, conservado en
\texttt{evidencias/muestra.sha256}, es:
\begin{comandos}
a1a442ca2e8062901ba3fc0dfd58f9f412e14c22332ee2e4698eeffab6bd787c
\end{comandos}
La prueba demuestra que funciona el análisis offline con esta muestra;
\textbf{no demuestra soporte monitor/inyección ni recuperación del router del aula}.

\subsection*{Verificación del ayudante}
\texttt{python3 -m unittest -v test\_lab10.py} finalizó con doce pruebas
aprobadas. Comprueban: rechazo de parámetros inválidos (radio, cantidad de ráfagas
y direcciones IPv4); simulación sin ejecutar herramientas ni crear salidas; captura
inexistente y conservación de clave existente; paso literal de rutas con espacios y
propagación de fallo; construcción correcta de los comandos de \texttt{monitor},
\texttt{escanear}, \texttt{inyectar}, \texttt{deauth} y \texttt{bloquear}; y
recuperación real de la clave pública. El registro es
\texttt{evidencias/pruebas\_locales.txt}. El diagnóstico local detectó la suite
e \texttt{iw}, pero ninguna interfaz Wi-Fi; se conserva su salida con código 1
en \texttt{evidencias/diagnostico.txt}.

\subsection*{Registro pendiente para clase}
Completar \path{evidencias/REGISTRO_CLASE.md} con BSSID, SSID y canal autorizados,
adaptador/controlador, inicio y fin, captura utilizada, IVs, RXQ, Lost y resultado
validado por el docente. Conservar pantalla de monitor, captura WEP y resultado
de recuperación, además del SHA-256 del archivo original.

Para comparar condiciones, tomar una línea base de 60 segundos, una ventana
controlada de igual duración solo si el docente la dispone y un retorno a la
línea base. Mantener constantes equipo, BSSID y canal. Calcular
\(r_{IV}=(\#Data_{final}-\#Data_{inicial})/\Delta t\).
Anotar N/D si falta una métrica; no sustituirla por cero ni inventar resultados.

\newpage
\section{Diagnóstico y conclusiones}
\begin{tabularx}{\linewidth}{|p{4cm}|X|}\hline
Síntoma & Revisión y siguiente paso \\\hline
No aparece Wi-Fi & Revisar acceso físico USB y controlador. En WSL2, Ethernet
virtual no basta; elegir un entorno con el adaptador expuesto. \\\hline
Monitor no compatible & Revisar modos de \texttt{iw list} y el chipset/controlador;
la instalación de la suite no modifica las capacidades de radio. \\\hline
Aviso de canal incorrecto & Verificar BSSID, canal, interfaz monitor y procesos
que cambien la sintonía; evitar otra captura saltando canales en el mismo radio. \\\hline
Beacons pero pocos IVs & Confirmar WEP, tráfico de datos y posición del receptor.
Un AP inactivo no proporciona automáticamente material suficiente. \\\hline
No se recupera la clave & Conservar captura completa, reunir más tráfico útil,
verificar BSSID y revisar ARP/PTW; no usar diccionarios como sustituto. \\\hline
Salida de clave ya existente & Usar otra ruta con \texttt{--clave}; el ayudante
rechaza sobrescribirla. Con comandos manuales, \texttt{-l} sí puede sobrescribir. \\\hline
\end{tabularx}

La recuperación de la muestra confirma experimentalmente la debilidad de WEP
sin diccionario. La ejecución en el aula dependerá del adaptador, el tráfico
disponible y el alcance fijado por el docente. Congestión y pérdida pueden
retrasar la obtención de IVs útiles; los contadores orientan la investigación,
pero no prueban por sí solos un ataque. Para redes de uso real se debe sustituir
WEP por un mecanismo de seguridad vigente.

\section*{Referencias}
\small
Enunciado local: \texttt{Laboratorio\_\_10.pdf}, una página. Fuentes primarias
consultadas el 20 de octubre de 2026:
\begin{enumerate}
\item Tews, Weinmann y Pyshkin. \textit{Breaking 104 bit WEP in less than
60 seconds}. IACR, 2007. \url{https://eprint.iacr.org/2007/120}
\item Aircrack-ng: recuperación WEP y PTW.
\url{https://aircrack-ng.org/doku.php?id=aircrack-ng}
\item Aircrack-ng: Airmon-ng.
\url{https://aircrack-ng.org/doku.php?id=airmon-ng}
\item Aircrack-ng: Airodump-ng y contadores.
\url{https://aircrack-ng.org/doku.php?id=airodump-ng}
\item Aircrack-ng: prueba de inyección.
\url{https://aircrack-ng.org/doku.php?id=injection_test}
\item Aircrack-ng: README y muestra pública.
\url{https://github.com/aircrack-ng/aircrack-ng/blob/master/README};
\url{https://dl.aircrack-ng.org/ptw.cap}
\item Microsoft Learn: conexión USB en WSL.
\url{https://learn.microsoft.com/en-us/windows/wsl/connect-usb}
\item Aircrack-ng: reinyección ARP.
\url{https://www.aircrack-ng.org/doku.php?id=arp-request_reinjection}
\item Wireshark: referencia de campos IEEE 802.11.
\url{https://www.wireshark.org/docs/dfref/w/wlan.html}
\end{enumerate}
\end{document}
